\section{Final-Checkpoint Family Fit: A Diagnostic Non-Result}
\label{sec:final}

Our \emph{first} attempt at validating Eq.~\eqref{eq:uet} was the
natural one: load the final-step checkpoints of all four Pythia-deduped
sizes we can run ($70$M, $160$M, $410$M, $1$B), evaluate on a fixed
test set, and fit $(c, \Linf)$ across the family. The fit returned
$c = 1.0$, $R^2 \approx -7 \times 10^{-8}$ --- numerically indistinguishable
from zero. This is not a bug in the fitter; it is a rank-deficient problem.
We document it here because the failure mode is instructive and because
the curriculum fit of \S\ref{sec:curriculum} is the principled response.

\subsection{Why the final-checkpoint fit is rank-deficient}

All four checkpoints were evaluated at $n = 287{,}744$ tokens of held-out
Pile; in the UET formula the feature $x = \deff \log(d/\deff)/n$ is
dominated by $1/n \approx 3.5 \cdot 10^{-6}$, the \emph{same} for every
model. With $x$ essentially constant across rows, the linear problem
\[
L_i \;=\; c \cdot x_i + \Linf + \varepsilon_i
\]
has condition number $\gg 10^{6}$ and $c$ is not identifiable. The
solver dutifully places all explanatory power in $\Linf$ and leaves
$c$ at its initialisation. The only correct response is to change the
experimental design --- which is what curriculum fitting accomplishes.

\subsection{What the final-checkpoint data \emph{does} show}

\begin{table}[H]
\centering
\small
\begin{tabular}{lSSSS}
\toprule
Model & {$d$ (hidden)} & {$N$ (params)} & {val loss} & {$\deff$} \\
\midrule
pythia-70m   & 512  & 70385664   & 4.160 & 132.07 \\
pythia-160m  & 768  & 162201600  & 3.555 & 45.12  \\
pythia-410m  & 1024 & 405012480  & 3.072 & 87.03  \\
pythia-1b    & 2048 & 1011351552 & 2.842 & 50.25  \\
\bottomrule
\end{tabular}
\caption{Final-step evaluation on 287{,}744 held-out tokens.
Loss decreases monotonically with model width, but $\deff$ is
\emph{non-monotone} because $\deff$ depends on spectral concentration,
not model capacity --- a larger model can distribute variance over more
directions (higher $\deff$) \emph{or} concentrate on fewer (lower
$\deff$), depending on the training objective and regularisation. This
is consistent with UET: $\deff$ is not a proxy for parameter count.}
\label{tab:scaling-final}
\end{table}

\begin{figure}[H]
\centering
\begin{tikzpicture}
\begin{axis}[
  width=0.80\textwidth, height=5.2cm,
  xmode=log, log basis x=10,
  xlabel={parameters $N$ (log)},
  ylabel={val loss},
  grid=both, grid style={line width=0.1pt, draw=gray!25},
  tick label style={font=\small}, label style={font=\small},
  legend style={font=\scriptsize, draw=none, at={(0.98,0.98)}, anchor=north east},
]
  \addplot+[mark=*, thick, color=blue!70!black, mark size=3pt]
    table[x=n_params, y=val_loss] {data/scaling_final.dat};
  \addlegendentry{Pythia family, final step}
\end{axis}
\end{tikzpicture}
\caption{Final-step loss vs. parameter count. A classical Chinchilla-style
power law fits this curve reasonably ($\alpha \approx 0.53$), but the
UET fit over the \emph{same points} fails because the design matrix is
degenerate in $n$.}
\label{fig:scaling-final}
\end{figure}

\begin{remark}[Why this matters for reviewers]
A reader might ask: ``If the final-checkpoint fit fails, how can you
claim $c$ is universal?'' The answer is that the final-checkpoint
setup cannot \emph{test} universality at all: it is blind to $n$. The
curriculum fit of \S\ref{sec:curriculum} varies $n$ by four orders of
magnitude within each model, which is the regime where Eq.~\eqref{eq:uet}
is empirically testable. The agreement of $c$ across 160M and 410M in
that regime is the load-bearing evidence.
\end{remark}
